\section{Preliminaries} \label{sec:prelim}

In this section we present the dif\/ferential geometric tools
we need to formulate and study the projective metrizability problem.

\looseness=1
A systems of second-order ordinary dif\/ferential equations on a
manifold $M$ can be identif\/ied with a second-order vector f\/ield
that is called a semispray. To each semispray one can associate a
geometric apparatus very useful to obtain qualitative information
regarding: the variations of its geodesics, their stability, as
well as the inverse problem of the calculus of variations,~\cite{bcd10}.
A global formulation for the geometric apparatus one can associate
to a semispray is due to Grifone \cite{grifone72} and can be
obtained using the Fr\"olicher--Nijenhuis theo\-ry~\cite{frolicher56}.

\subsection[Fr\"olicher-Nijenhuis theory]{Fr\"olicher--Nijenhuis theory} \label{subsec:fn}

In this subsection we recall and extend some aspects of the
Fr\"olicher--Nijenhuis theory, which will be applied in the next
subsection to vector valued dif\/ferential forms on tangent bundles.
For the classic and modern formulations of Fr\"olicher--Nijenhuis theory we
refer to \cite{frolicher56, grifone72, grifone00, KMS93, deleon89}.

In this work $M$ is a real, $n$-dimensional, smooth manifold. We
denote by $C^{\infty}(M)$, the ring of smooth functions on $M$,
and by ${\mathfrak X}(M)$, the $C^{\infty}(M)$-module of vector
f\/ields on $M$. Consider $\Lambda(M)=\bigoplus_{k\in {\mathbb
N}}\Lambda^k(M)$ the graded algebra of dif\/ferential forms on $M$. We
denote by $S^k(M)$ the space of symmetric $(0,k)$ tensors on $M$.
We also write $\Psi(M)=\bigoplus_{k\in {\mathbb N}}\Psi^k(M)$ for the
graded algebra of vector-valued dif\/ferential forms on $M$.

For $L\in \Psi^l(M)$, a vector valued $l$-form, we consider
$\tau_L: \Lambda^1(M) \otimes \Lambda^k(M) \to \Lambda^{k+l}(M)$,
or $\tau_L: \Psi^1(M) \otimes \Psi^k(M) \to \Psi^{k+l}(M)$, the
\emph{alternating operator} def\/ined as
\begin{gather}
\label{taull}  (\tau_L B)(X_1,\dots ,X_{k+l}) = \frac{1}{k!l!} \sum_{\sigma \in S_{k+l}}
\varepsilon(\sigma) B(L(X_{\sigma(1)},\dots , X_{\sigma(l)}),
X_{\sigma(l+1)},\dots , X_{\sigma(l+k)}),
\end{gather}
where
$X_1,\dots , X_{k+l} \in {\mathfrak X}(M)$ and $S_{k+l}$ is the
permutation group of $\{1,\dots , k+l\}$.

The restriction of $\tau_L$ to $\Lambda^{k+1}(M)\subset
\Lambda^1(M) \otimes \Lambda^k(M)$, or the restriction to
$\Psi^{k+1}(M)$, is a~derivation of degree $(l-1)$ and it
coincides with the \emph{inner product}~$i_L$, see \cite{grifone00,
KMS93}.  Inner product~$i_L$
is trivial on $\Lambda^0(M)=C^{\infty}(M)$, or
$\Psi^0(M)=\mathfrak{X}(M)$, and hence it is a \emph{derivation of type
$i_*$} \cite{grifone00} or an \emph{algebraic derivation}
\cite{KMS93}. Since it satisf\/ies the Leibniz rule, $i_L$ is uniquely
determined by its action on $\Lambda^1(M)$, or $\Psi^1(M)$, when it is
given by $i_LB=B\circ L$. For the particular case when $l=1$ and $L=\operatorname{Id}$ we have that
$i_{\operatorname{Id}}B=kB$ for all $B\in \Lambda^k(M)$, or $B\in
\Psi^k(M)$.

For a linear connection $\nabla$ on $M$ consider
$d^{\nabla}:\Psi^k(M)\to \Psi^{k+1}(M)$ the covariant exterior
derivative, see \cite[\S~11.13]{KMS93}, given by
\begin{gather} \nonumber   d^{\nabla}B(X_1,\dots , X_{k+1}) =
\sum_{i=1}^{k+1}
(-1)^{i+1}\nabla_{X_i}B(X_1,\dots , \hat{X}_i, \dots , X_{l+1})  \\
\label{dnablaL}
\phantom{d^{\nabla}B(X_1,\dots , X_{k+1}) =}{}
+ \sum_{1\leq i<j\leq k+1} (-1)^{i+j}B([X_i,
X_j], X_1,\dots , \hat{X}_i,\dots , \hat{X}_j,\dots , X_{k+1}).\!\!\!
\end{gather} The exterior derivative $d: \Lambda^k(M)\to
\Lambda^{k+1}(M)$ satisf\/ies also formula \eqref{dnablaL} for $B\in
\Lambda^k(M)$. Therefore, we will use the notation $d^{\nabla}$ to refer to
both, the covariant exterior derivative, or the exterior
derivative. For the latter case $d=d^{\nabla}$ does not depend on
the linear connection $\nabla$.

For a vector valued $l$-form $L$, consider the commutator of the inner
product $i_L$ and the (covariant) exterior derivative $d$
($d^{\nabla}$). This dif\/ferential operator is
denoted by  $d_L: \Lambda^k(M) \to \Lambda^{k+l}(M)$ ($d^{\nabla}_L:
\Psi^k(M) \to \Psi^{k+l}(M)$), it is given by
\begin{gather} d^{\nabla}_L=i_L\circ d^{\nabla}+(-1)^{l}d^{\nabla}\circ i_L,
\label{dil}
\end{gather} it is a derivation of degree $l$, which is
called the \emph{$($covariant$)$ exterior derivative with respect to}~$L$. Derivation $d_L$ ($d^{\nabla}_L$) commutes with the exterior
derivative $d$ ($d^{\nabla}$) and hence it is a \emph{derivation of type}
$d_*$ \cite{grifone00} or a \emph{Lie derivation} \cite{KMS93}.
Since it satisf\/ies the Leibniz rule, $d_L$ ($d_{L}^{\nabla}$) is
uniquely determined by its action on $\Lambda^0(M)=C^{\infty}(M)$
($\Psi^0(M)={\mathfrak X}(M)$). For the particular case when $l=1$ and
$L=\operatorname{Id}$ we have that
$d^{\nabla}_{\operatorname{Id}}=d^{\nabla}$. Therefore, we obtain
$d^{\nabla}\operatorname{Id}=T$, where $T$ is the torsion of the linear
connection $\nabla$.

For two vector valued forms $L\in \Psi^l(M)$ and $K\in \Psi^k(M)$,
their Fr\"olicher--Nijenhuis bracket $[L,K]$ is a vector
valued $(k+l)$-form, def\/ined by
\begin{gather} \label{dlk}
d_{[L,K]}=d_L\circ d_K -(-1)^{kl}d_K\circ d_L.
\end{gather}

For a vector valued $l$-form $L$ and a linear connection $\nabla$
on $M$ we obtain a derivation of degree~$l$ given by ${\mathcal
D}_L=\tau_L \nabla$. Hence, ${\mathcal D}_L: \Lambda^k(M) \to
\Lambda^{k+l}(M)$ (${\mathcal D}_L: \Psi^k(M) \to
\Psi^{k+l}(M)$), acts on (vector-valued) $k$-forms as follows:
\begin{gather} \left({\mathcal
D_L}B\right)(X_1,\dots ,X_{k+l})  = \frac{1}{l!k!}
\sum_{\sigma\in S_{k+l}}
\varepsilon(\sigma)\big(\nabla_{L(X_{\sigma(1)},\dots ,X_{\sigma(l)})}B\big)\left(X_{\sigma(l+1)},\dots ,
X_{\sigma(k+l)}\right). \label{DLlk}
\end{gather}

For the particular case when $l=1$ and $L=\operatorname{Id}$, we
denote the corresponding derivation of degree $1$ by $\mathcal{D} =
\mathcal{D}_{\operatorname{Id}}$.
Since any derivation of degree $l$ can be uniquely
decomposed into a sum of a Lie derivation and an algebraic
derivation \cite[\S~8.3]{KMS93}, we obtain for ${\mathcal D}_L$
the following result.

\begin{Lemma} \label{lem:dec_DL} For a vector valued $l$-form $L$ and a linear connection $\nabla$ on $M$,
derivation of degree~$l$,~${\mathcal D}_L$, decomposes uniquely
into a sum of a Lie derivation and an algebraic derivation
as follows:
\begin{gather}
{\mathcal D}_L= d^{\nabla}_L -
i_{d^{\nabla}_L\operatorname{Id}}.\label{decompDL}
\end{gather}
\end{Lemma}
\begin{proof}
When acting on forms, formula \eqref{decompDL} reads ${\mathcal
D}_L= d_L - i_{d^{\nabla}_L\operatorname{Id}}$. The vector valued
$(l+1)$-form that def\/ines the inner product in formula
\eqref{decompDL} is, according to formula \eqref{dil}, given by
$d^{\nabla}_L\operatorname{Id}=i_LT +(-1)^ld^{\nabla}L.$

Since Lie derivations commute with the exterior
derivative $d^{\nabla}$ and satisfy the Leibnitz rule it follows that
they are uniquely determined by their action on
$\Lambda^0(M)=C^{\infty}(M)$ ($\Psi^0(M)={\mathfrak X}(M)$).
Using formulae \eqref{dil} and \eqref{DLlk} one can immediately
check that ${\mathcal D}_Lf=d_Lf$ for any scalar (vector
valued) $0$-form $f$.

Since algebraic derivations are trivial on
$\Lambda^0(M)=C^{\infty}(M)$ ($\Psi^0(M)={\mathfrak X}(M)$), and
satisfy the Leibnitz rule it follows that they are uniquely determined by
their action on $\Lambda^1(M)$ ($\Psi^1(M)$). To prove formula
\eqref{decompDL} we have to show that
\[
\left({\mathcal D}_L- d^{\nabla}_L\right)\omega = -\omega \circ
\left(d^{\nabla}_L\operatorname{Id}\right),
\]
 for any (vector valued) $1$-form
$\omega$. Since formally, we have the same formulae to def\/ine the
action on scalar, or vector valued forms, we will work with scalar
forms.

Let $\omega \in \Lambda^1(M)$ and $X_1, \dots , X_{l+1} \in
{\mathfrak X}(M)$. For $k=1$, from formula \eqref{DLlk} we obtain
the action of the derivation ${\mathcal D}_L$ on $1$-forms as
follows:
\begin{gather}
 \left({\mathcal D}_L\omega\right)(X_1,\dots , X_{l+1})=
\sum_{i=1}^{l+1}(-1)^{l+1-i}\left(\nabla_{L(X_1,\dots ,
\hat{X}_i,\dots ,X_{l+1})}\omega\right)(X_i) \nonumber  \\
\qquad{}
=  \sum_{i=1}^{l+1}(-1)^{l+1-i} \left\{L(X_1,\dots , \hat{X}_i,\dots ,X_{l+1})
(\omega(X_i)) -\omega\left( \nabla_{L(X_1,\dots ,
\hat{X}_i,\dots ,X_{l+1})} X_i\right)\right\}. \!\!\!\label{eq:dlo1}
\end{gather}

From formula \eqref{dil} we obtain that the action of the exterior
derivative $d_L$ on a $1$-form $\omega$ is given by
$d_L\omega=i_Ld\omega+ (-1)^ld(\omega\circ L)$.
Therefore, for $X_1, \dots , X_{l+1} \in {\mathfrak X}(M)$ we have{\samepage
\begin{gather}
\nonumber (d_L\omega)(X_1,\dots ,X_{l+1})   =   \sum_{i=1}^{l+1}(-1)^{l+1-i}
L(X_1,\dots , \hat{X}_i,\dots , X_{l+1})(\omega(X_i))  \\
\label{eq:dlo2}
\phantom{(d_L\omega)(X_1,\dots ,X_{l+1})   =}{}
+\sum_{i=1}^{l+1}(-1)^{l+1-i} \omega([X_i, L(X_1,\dots ,
\hat{X}_i,\dots , X_{l+1})])  \\
\nonumber   \phantom{(d_L\omega)(X_1,\dots ,X_{l+1})   =}{}
 +\sum_{1\leq i<j\leq l+1}
(-1)^{l+i+j}\omega(L([X_i,X_j], X_1,\dots , \hat{X}_i,\dots ,
\hat{X}_j,\dots , X_{l+1})).
\end{gather}
Now, we evaluate $d^{\nabla}_L\operatorname{Id}=i_LT + (-1)^l
d^{\nabla}L$ on $l+1$ vectors $X_1, \dots , X_{l+1} \in {\mathfrak
X}(M)$.}

For $k=1$, if we restrict the action of $\tau_L$ given by formula
\eqref{taull} to $\Psi^2(M)$ we obtain that the inner product $i_L: \Psi^2(M) \to
\Psi^{l+1}(M)$ is given by:
\begin{gather}
\left(i_LT\right)(X_1,\dots ,
X_{l+1})=\sum_{i=1}^{l+1}(-1)^{l-i}T(X_i, L(X_1,\dots , \hat{X}_i,
\dots , X_{l+1})). \label{eq:ilt}
\end{gather} Using formula \eqref{dnablaL}, the action of the exterior covariant derivative
$d^{\nabla}$ on the vector valued
$l$-form~$L$ is given by
\begin{gather}
\nonumber   d^{\nabla}L(X_1,\dots , X_{l+1}) = \sum_{i=1}^{l+1}
(-1)^{i+1}\nabla_{X_i}L(X_1,\dots , \hat{X}_i, \dots , X_{l+1})  \\
\phantom{d^{\nabla}L(X_1,\dots , X_{l+1}) =}{}
 + \sum_{1\leq i<j\leq l+1} (-1)^{i+j}L([X_i, X_j], X_1,\dots ,
\hat{X}_i,\dots , \hat{X}_j,\dots , X_{l+1}). \label{eq:dnl}
\end{gather}
Using formulae \eqref{eq:dlo1}, \eqref{eq:dlo2}, \eqref{eq:ilt}, and \eqref{eq:dnl} it follows that
\[
\left({\mathcal D}_L \omega-d_L\omega\right)(X_1,\dots ,X_{l+1}) =
-\left(\omega\circ (i_LT+(-1)^ld^{\nabla}L)\right)(X_1,\dots ,X_{l+1}),
\]
for all $X_1,\dots , X_{l+1}\in  {\mathfrak X}(M)$, which means that the decomposition \eqref{decompDL} is true.
\end{proof}

For the particular case when $l=1$ and $L=\operatorname{Id}$, we have that
the vector valued $2$-form
$d_{\operatorname{Id}}^{\nabla}\operatorname{Id}$ reduces to
torsion $T$ since $i_{\operatorname{Id}}T=2T$,
$d^{\nabla}\operatorname{Id}=T$. Therefore, decomposition
\eqref{decompDL} becomes
\[ %\label{DdiT}
{\mathcal D}=d- i_T.
\]
\begin{Remark}
Formula \eqref{decompDL} shows that the dif\/ference of the two
derivations $d_L-{\mathcal
D}_L=i_{d^{\nabla}_L\operatorname{Id}}$ is an algebraic derivation.
In other words, if $\omega \in \Lambda^k(M)$ vanishes at some
point $p\in M$, $\omega_p=0$, then $({\mathcal
D}_L\omega)_p=(d_L\omega)_p$. For the particular case when $l=1$, this result has been
shown in \cite[Proposition~2.5]{grifone00}.
\end{Remark}



\subsection[Vertical calculus on $TM$ and semi-basic forms]{Vertical calculus on $\boldsymbol{TM}$ and semi-basic forms} \label{subsec:vcalc}

Consider $(TM, \pi, M)$, the tangent bundle of the manifold $M$
and $(TM\setminus\{0\}), \pi, M)$ the slashed tangent bundle, which is
the tangent bundle with the zero
section removed. The tangent bundle carries some canonical structures, such
as the vertical distribution, the Liouville vector f\/ield, and the
vertical endomorphism. The dif\/ferential calculus associated to
these structures, using the Fr\"olicher--Nijenhuis theory developed in
the previous subsection, plays an important role in the geometry of a
system of second-order ordinary dif\/ferential equations,
\cite{ bcd10, bucataru09, grifone72, grifone00, klein68, deleon89, miron94}.

The \emph{vertical subbundle} is def\/ined as $VTM=\{\xi \in TTM,
(D\pi)(\xi)=0\}$. It induces a vertical distribution $V: u\in TM
\mapsto V_u=VTM\cap T_uTM$. This distribution is $n$-dimensional
and it is integrable, being tangent to the leaves of the \emph{natural
foliation} induced by submersion~$\pi$. If~$(x^i)$ are local coordinates on the base
manifold~$M$, we denote by $(x^i, y^i)$ the induced coordinates on~$TM$. It follows that~$y^i$ are coordinates along the leaves of the
natural foliation, while $x^i$ are transverse coordinates for the
foliation. We denote by ${\mathfrak X}^v(TM)$ the Lie subalgebra of vertical vector f\/ields on~$TM$. An
important vertical vector f\/ield on $TM$ is the \emph{Liouville vector
field}, which locally is given by
$\mathbb{C}=y^i{\partial}/{\partial y^i}.$

The \emph{tangent structure} (or vertical endomorphism) is the
$(1,1)$-type tensor f\/ield $J$ on~$TM$, which locally can be
written as follows:
\[
J=\frac{\partial}{\partial y^i} \otimes dx^i.
\]
Tensor $J$ satisf\/ies $J^2=0$ and $\operatorname{Ker} J=
\operatorname{Im} J=VTM$. Tangent structure $J$ is an integrable structure
since the Fr\"olicher--Nijenhuis bracket vanishes, $[J,J]=0$. As a consequence and using formula~\eqref{dlk} we have that
$d_J^2=0$.

For the natural foliation induced by submersion~$\pi$ and the
corresponding vertical distribution there are some important classes
of forms: basic and semi-basic forms. As we will see in the next
sections, semi-basic forms, vector valued semi-basic forms, and vector valued almost
semi-basic forms are important ingredients to formulate and
address the projective metrizability problem.
\begin{Definition} Consider $\omega\in \Lambda(TM)$ and $L\in \Psi(TM)$.
\begin{itemize}\itemsep=0pt
\item[i)] $\omega$ is called a \emph{basic} form if both $\omega$ and
  $d\omega$ vanish whenever one of the arguments of $\omega$
  (respectively $d\omega$) is a vertical vector f\/ield.
\item[ii)] $\omega$ is called a \emph{semi-basic} form if it vanishes whenever
one of its arguments is a vertical vector f\/ield.
\item[iii)] $L$ is called a vector valued \emph{semi-basic} form if it takes vertical values and vanishes
whenever one of its arguments is a vertical vector f\/ield.
\item[iv)] $L$ is called a vector valued \emph{almost semi-basic} form
  if it vanishes whenever one of its arguments is a vertical vector f\/ield
and for every vertical vector f\/ield $X\in {\mathfrak X}^v(TM)$ we
have that $\mathcal{L}_XL$ is a vector valued semi-basic form.
\end{itemize}
\end{Definition}
In local coordinates a basic $k$-form $\omega$ on $TM$ can
be written as
\begin{gather*}
 \omega=\frac{1}{k!}\omega_{i_1\dots i_k}(x)dx^{i_1}\wedge \cdots
\wedge dx^{i_k}.%\label{eq:bf}
\end{gather*}
For basic forms, the coordinates functions $\omega_{i_1\dots i_k}(x)$ are basic
functions, which means that they are constant along the leaves of the
natural foliation.

Locally, a semi-basic $k$-form $\omega$ on $TM$ can
be written as
\begin{gather}
 \omega=\frac{1}{k!}\omega_{i_1\dots i_k}(x,y)dx^{i_1}\wedge \cdots
\wedge dx^{i_k}.\label{eq:smf}
\end{gather}
We will denote by $\Lambda^k_v(TM)$ the set of
semi-basic $k$-forms on $TM$. A $1$-form $\omega$ on $TM$ is
semi-basic if and only if $i_J\omega=\omega \circ J=0$.

In local coordinates, a vector valued semi-basic $l$-form $L$ on
$TM$ can be written as
\begin{gather}
L=\frac{1}{l!}L^{\underline{j}}_{i_1\dots i_l}(x,y)\frac{\partial}{\partial y^j}\otimes dx^{i_1}\wedge \cdots
\wedge dx^{i_l}.\label{eq:vsmf}
\end{gather} In this work all contravariant or covariant
indices, of some tensorial coef\/f\/icients, that refer to vertical components will be underlined. We
will denote by $\Psi^l_v(TM)$ the set of vector valued semi-basic
$l$-forms on $TM$. A vector valued $1$-form $L$ on $TM$ is
semi-basic if and only if $J\circ L=0$ and $i_JL=L \circ J=0$. The
tangent structure $J$ is a vector valued semi-basic $1$-form.

Locally, a vector valued almost semi-basic $l$-form $L$ on $TM$ can be
expressed as
\begin{gather}
 L =\frac{1}{l!}L^{j}_{i_1\dots i_l}(x)\frac{\partial}{\partial
x^j}\otimes dx^{i_1}\wedge \cdots \wedge dx^{i_l}  +
\frac{1}{l!}L^{\underline{j}}_{i_1\dots i_l}(x,y)\frac{\partial}{\partial
y^j}\otimes dx^{i_1}\wedge \cdots \wedge
dx^{i_l}.\label{eq:vasmf}
\end{gather}
For a vector $X$ on $TM$ and a vector valued $l$-form $L$
on $TM$, the Fr\"olicher--Nijenhuis bracket $[X,L]$, def\/ined by formula
\eqref{dlk}, is a vector valued $l$-form on $TM$ given by:
\begin{gather*}
[X,L](X_1,\dots ,X_l)=[X,L(X_1,\dots ,X_l)] - \sum_{i=1}^lL([X,X_i],
X_1,\dots ,\hat{X}_i,\dots ,X_l),  %\label{eq:xl}
\end{gather*}
for $X_1, \dots , X_l$ vector f\/ields on $TM$. Using the above formula and the fact that the vertical distribution is integrable it
follows that vector valued semi-basic forms are also almost
semi-basic. This can be seen also from the local expressions
\eqref{eq:vsmf} and \eqref{eq:vasmf}.

Next two lemmas give a good motivation for considering the class of vector
valued almost semi-basic forms. We will also see in Section~\ref{sec:fi} that the partial
dif\/ferential operator we use to discuss the projective metrizability
problem is def\/ined in terms of some vector valued almost
semi-basic forms.

\begin{Lemma} \label{lem:dlomega} Let $L$ be a vector valued almost semi-basic $l$-form on~$TM$.
Then, the differential operator $d_L$ preserves semi-basic forms,
$d_L: \Lambda^k_v(TM) \to \Lambda^{k+l}_v(TM).$
\end{Lemma}

\begin{proof}
Consider a vector valued almost semi-basic $l$-form $L$, locally given
by formula \eqref{eq:vasmf}, and a semi-basic $k$-form $\omega$, locally given by formula
\eqref{eq:smf}. Since $d_L$ is a derivation of degree $l$, it follows
that the $(k+l)$-form $d_L\omega$ can be
expressed locally as follows
\begin{gather}
\nonumber d_L\omega  =  \frac{1}{k!}d_L(\omega_{i_1\cdots
  i_k})\wedge  dx^{i_1}\wedge \cdots \wedge dx^{i_k}  \\[-0.5ex]
 \phantom{d_L\omega  =}{}  +  \frac{1}{k!}
\sum^k_{j=1}(-1)^{(j-1)l} \omega_{i_1\cdots
  i_k} dx^{i_1}\wedge \cdots \wedge d_Ldx^{i_j} \wedge \cdots \wedge
dx^{i_k}. \label{eq:dlo}
\end{gather}
Using the assumption that $L$ is a vector valued almost semi-basic
form, we show that all terms in the right hand side of the above formula
are semi-basic forms. Since $L$ vanishes whenever one of its arguments
is a vertical vector f\/ield it follows that for a function $f\in C^{\infty}(TM)$, $d_Lf=i_Ldf=df\circ
L$ is a semi-basic $l$-form. Hence $d_L(\omega_{i_1\cdots i_k})$ are semi-basic $l$-forms.

We will prove now that $d_Ldx^{i_j}=(-1)^ldd_Lx^{i_j} $ are semi-basic
$(l+1)$-forms. Using the local expression \eqref{eq:vasmf} of $L$, we have
\begin{gather*} d_Lx^{i_j}=i_Ldx^{i_j} = dx^{i_j}\circ L =\frac{1}{l!}
L^{i_j}_{i_1\dots i_l}(x) dx^{i_1}\wedge \cdots \wedge
dx^{i_l},%\label{eq:dlx}
\end{gather*}
which are basic $l$-forms. Therefore, $dd_Lx^{i_j}$ are basic and hence
semi-basic $(l+1)$-forms.
One can conclude now that all terms in the right hand side of formula~\eqref{eq:dlo} are
semi-basic forms and hence~$d_L\omega$ is a semi-basic~$(k+l)$-form.
\end{proof}

\begin{Lemma} \label{lem:nablaomega} Consider $\nabla$ a linear connection on $TM$ such
that $\nabla J=0$ and a vector valued almost semi-basic $l$-form
$L$ on $TM$. Then, the differential operator ${\mathcal D}_L$
preserves semi-basic forms, ${\mathcal D}_L: \Lambda^k_v(TM) \to
\Lambda^{k+l}_v(TM).$
\end{Lemma}

\begin{proof}
Using formula \eqref{decompDL} and Lemma \ref{lem:dlomega} we have
that the dif\/ferential operator ${\mathcal D}_L$ preserves semi-basic forms
if and only if the algebraic derivation of degree $l$,
$i_{d^{\nabla}_L\operatorname{Id}}$ preserves semi-basic forms. Hence,
we will complete the proof if we show that the vector valued
$(l+1)$-form $d^{\nabla}_L\operatorname{Id}=i_LT+(-1)^ld^{\nabla}_L$
takes vertical values whenever one of its arguments is a vertical
vector f\/ield. Here~$T$ is the torsion of the linear connection $\nabla$.

Using formulae \eqref{eq:ilt} and \eqref{eq:dnl} we have
\begin{gather}
\nonumber   (-1)^l(d^{\nabla}_L\operatorname{Id})(X_1,\dots ,X_{l+1}) =(-1)^l(i_LT+(-1)^ld^{\nabla}_L)(X_1,\dots ,X_{l+1})
\\[-0.5ex] \nonumber
\qquad{}  =  \sum_{i=1}^{l+1}(-1)^{i+1}
\left\{\nabla_{L(X_1,\dots ,\hat{X}_i,\dots ,X_{l+1})}X_i + [X_i,
  L(X_1,\dots ,\hat{X}_i,\dots ,X_{l+1})]\right\} \\[-0.5ex]
 \qquad\quad{}  +  \sum_{1\leq i<j\leq
  l+1} (-1)^{i+j}L([X_i,X_j], X_1,\dots , \hat{X}_i,\dots , \hat{X}_j,\dots ,
X_{l+1}). \label{eq:dnlid}
\end{gather}
We will show now that whenever one of the
$(l+1)$ arguments of $d^{\nabla}_L\operatorname{Id}$ is a vertical
vector f\/ield, then the right hand side of formula~\eqref{eq:dnlid} is
a vertical vector f\/ield. Using the fact that $d^{\nabla}_L\operatorname{Id}$ is a~vector valued $(l+1)$-form, we will discuss only the case when $X_1$ is a vertical
vector f\/ield. Since~$L$ vanishes whenever one of its arguments is a
vertical vector f\/ield, the nonzero vector f\/ield that remains from the
right hand side of formula~\eqref{eq:dnlid}, when $X_1$ is a vertical
vector f\/ield, is
\begin{gather}
\nabla_{L(X_2,\dots ,X_{l+1})}X_1+[X_1,L](X_2,\dots ,X_{l+1}). \label{eq:nl1}
\end{gather} The condition $\nabla J=0$ implies that the linear
connection $\nabla$ preserves the vertical distribution and since $X_1$ is a vertical
vector f\/ield it follows that $\nabla_{L(X_2,\dots,X_{l+1})}X_1 $ is a
vertical vector f\/ield as well. Since $L$ is a vector valued almost
semi-basic $l$-form and $X_1$ is a vertical
vector f\/ield it follows that $[X_1, L]$ is a vector valued semi-basic
form, therefore it takes values into the vertical distribution and
hence  $[X_1,L](X_2,\dots ,X_{l+1})$ is a vertical vector f\/ield. It
follows that the vector f\/ield in formula~\eqref{eq:nl1} is vertical
and hence we have completed the proof.
\end{proof}


\section{Projective metrizability problem of a spray}
\label{sec:pms}

A system of homogeneous second-order ordinary dif\/ferential equations
on a manifold~$M$, whose coef\/f\/icients do not depend explicitly on time, can be identif\/ied
with a special vector f\/ield on~$TM$ that is called a spray. In
this section we address the following question, known as the
projective metrizability problem: for a given spray $S$ f\/ind necessary
and suf\/f\/icient conditions for the existence of a Finsler function $F$
such that the geodesics of $S$ and the geodesics of $F$ coincide up to
an orientation preserving reparameterization. We obtain such necessary and
suf\/f\/icient conditions in Theorem \ref{thm:pm} and these conditions are expressed
in terms of semi-basic $1$-forms.

Particular aspects of the projective metrizability problem were studied more than a century ago by Hamel
\cite{hamel03}. The problem was formulated rigorously in 1960's by Rapcs\'ak
\cite{rapcsak62} and Klein and Voutier \cite{klein68}. Yet, the projective
metrizability problem is far from being solved, and in the last decade
it has been intensively studied \cite{paiva05, bucataru09,
  crampin07b, crampin08,  shen01, szilasi07,  szilasi02, yang11}.


\subsection{Spray, nonlinear connection, and curvature}
\label{subsec:sn}

In this subsection, we start with a spray $S$ and use the Fr\"olicher--Nijenhuis theory
to derive a dif\/ferential calculus on $TM\setminus\{0\}$ \cite{grifone72} and to
obtain information about the given system of SODE. For the remaining
part of the paper, all geometric objects will be considered def\/ined on
the slashed tangent bundle $TM\setminus\{0\}$ and not on the whole
$TM$. This is motivated by the fact that we will want to connect them
with geometric structures in Finsler geometry, where the Finsler
function is not dif\/ferentiable on the zero section.

\begin{Definition} \label{def:spray} A vector f\/ield $S\in{\mathfrak X}(TM\setminus\{0\})$
is called a \emph{spray} if
\begin{itemize}\itemsep=0pt
\item[i)] $JS=\mathbb{C}$,
\item[ii)] $[\mathbb{C}, S]=S$. \end{itemize}
\end{Definition}
First condition in Def\/inition~\ref{def:spray} expresses that a
spray $S$ can be locally given as
\[
 S=y^i\frac{\partial}{\partial x^i} - 2G^i(x,y) \frac{\partial}{\partial   y^i},
  \]
   for some functions $G^i$ def\/ined on domains of induced
coordinates on $TM\setminus\{0\}$.

Second condition in Def\/inition~\ref{def:spray} expresses that the
vector f\/ield $S$ is $2$-homogeneous. It is equivalent with the
fact that functions $G^i$ are $2$-homogeneous in the f\/ibre
coordinates. In this work we will consider positive homogeneity
only and hence $G^i(x,\lambda y)=\lambda^2 G^i(x,y)$
for all $\lambda>0$. By Euler's theorem this homogeneity condition is equivalent to
$\mathbb{C}(G^i)=2G^i$.

A curve $c:I \to M$ is called \emph{regular} if its tangent lift takes
values in the slashed tangent bundle, $c': I \to
TM\setminus\{0\}$. A regular curve is called a \emph{geodesic} of
spray $S$ if $S\circ c'=c''$. Locally, $c(t)=(x^i(t))$ is a
geodesic of spray $S$ if
\begin{gather*}
\frac{d^2x^i}{dt^2}+2G^i\left(x,\frac{dx}{dt}\right)=0.
%\label{d2g}
\end{gather*}

\begin{Definition} A \emph{nonlinear connection} (or a horizontal
distribution, or Ehresmann connection) is def\/ined by an
$n$-dimensional distribution $H: u\in TM\setminus\{0\} \to
H_u\subset T_u(TM\setminus\{0\})$ that is supplementary to the
vertical distribution. \end{Definition}

Every spray induces a nonlinear connection through the
corresponding horizontal and vertical projectors, \cite{grifone72}
\begin{gather*}
h=\frac{1}{2}\left(\operatorname{Id}-\mathcal{L}_SJ\right), \qquad
v=\frac{1}{2}\left(\operatorname{Id}+\mathcal{L}_SJ\right).
%\label{hv}
\end{gather*} Locally, the two projectors $h$ and $v$ can
be expressed as follows
\[
h=\frac{\delta}{\delta x^i} \otimes dx^i, \qquad
v=\frac{\partial}{\partial y^i}\otimes \delta y^i,
\] where
\[
\frac{\delta}{\delta x^i}=\frac{\partial}{\partial
x^i}-N^j_i(x,y)\frac{\partial}{\partial y^j}, \qquad \delta y^i=dy^i+
N^i_j(x,y)dx^j, \qquad N^i_j(x,y)=\frac{\partial G^i}{\partial y^j}(x,y).
\]
Horizontal projector $h$ is a vector valued almost semi-basic
$1$-form.

For a spray $S$ consider the vector valued
semi-basic $1$-form
\begin{gather*}
\Phi=-v\circ \mathcal{L}_Sv = v\circ \mathcal{L}_Sh=v\circ
\mathcal{L}_S \circ h,
\end{gather*} which will be called the
\emph{Jacobi endomorphism}. It is also known as the Douglas tensor
\cite[Def\/inition~3.17]{grifone00} or as the Riemann curvature \cite[Def\/inition~8.1.2]{shen01}.
 Locally, the Jacobi endomorphism can be expressed as follows
\begin{gather*}
\Phi= R^i_j(x,y) \frac{\partial}{\partial y^i}
\otimes dx^j, \qquad R^i_j = 2\frac{\delta G^i}{\delta x^j} -
S(N^i_j) + N^i_kN^k_j. %\label{jacobi_end}
\end{gather*}

Another important geometric structure induced by a spray $S$ is
the \emph{curvature tensor} $R$. It is the vector valued semi-basic $2$-form
\begin{gather} R=\frac{1}{2}[h,h]=\frac{1}{2}R^i_{jk}\frac{\partial}{\partial y^i}\otimes dx^j\wedge dx^k. \label{curvature}
\end{gather}
Locally, the components of the curvature tensor, $R^i_{jk}$, are given by
\[
R^i_{jk}=\frac{\delta N^i_j}{\delta x^k} - \frac{\delta N^i_k}{\delta
x^j}.
\] Curvature tensor $R$ expresses the obstruction to the
integrability of the nonlinear connection.  Using formulae~\eqref{dlk}
and~\eqref{curvature} we have that $d_h^2=d_R$.

All the geometric objects induced by a spray $S$ inherit the
homogeneity condition. Therefore $[\mathbb{C}, h]=0$, which means
that the nonlinear connection is $1$-homogeneous. Also $[\mathbb{C},
R]=0$, $[\mathbb{C},\Phi]=\Phi$ and hence the the curvature tensor $R$
is $1$-homogeneous, while the  Jacobi endomorphism $\Phi$
is $2$-homogeneous.

Using the Jacobi identity, \cite[Proposition~2.7]{grifone00} , for the vector valued $0$-form $S$ and the
vector valued $1$-form $J$ we have $[J,[S,J]]-[J,[J,S]]-[S,[J,J]]=0$.
Therefore, we obtain $[J,h]=-2[J,[S,J]]=0$.

The two semi-basic vector vector valued $1$ and $2$-forms $\Phi$ and
$R$ are related as follows:
\begin{gather}
\Phi=i_SR, \qquad [J,\Phi]=3R. \label{eq:phir} \end{gather}
First formula in \eqref{eq:phir} is a consequence of the homogeneity,
while the second one is true in a~more general context. Locally, the
above two formulae can be expressed as follows:
\begin{gather*}
R^i_j=R^i_{kj}y^k, \qquad R^i_{jk}=\frac{1}{3}\left(\frac{\partial
    R^i_k}{\partial y^j}-\frac{\partial R^i_j}{\partial
    y^k}\right).%\label{eq:rr}
    \end{gather*}
An important class of sprays, which we will use in the last section to
provide examples of projectively metrizable sprays, is that of
isotropic sprays, \cite[Def\/inition~3.29]{grifone00}.
\begin{Definition} A spray $S$ is called \emph{isotropic} if its Jacobi
  endomorphism has the form
\begin{gather}
\Phi=\lambda J + \eta \otimes \mathbb{C}, \label{eq:phi_iso}
\end{gather}
where $\lambda\in C^{\infty}(TM\setminus\{0\})$ and $\eta$ is a
semi-basic $1$-form on $TM\setminus\{0\}$.
\end{Definition}
Due to f\/irst formula in \eqref{eq:phir} we have that $i_S\Phi=0$ and
hence $\lambda=-i_S\eta$. Also formulae \eqref{eq:phir} allows us to express
the isotropy condition \eqref{eq:phi_iso} for a spray in terms of the curvature tensor
$R$.
\begin{Proposition} \label{prop:iso}
A spray $S$ is isotropic if and only if its curvature tensor $R$ has the form
\begin{gather}
R=\alpha\wedge J + \beta \otimes \mathbb{C}, \label{eq:r_iso}
\end{gather}
where $\alpha$ is a semi-basic $1$-form and $\beta$ is a semi-basic
$2$-form on $TM\setminus\{0\}$. \end{Proposition}
\begin{proof}
We will prove that formulae \eqref{eq:phi_iso} and \eqref{eq:r_iso}
are equivalent.

Suppose that spray $S$ is isotropic. Therefore, the Jacobi endomorphism
$\Phi$ satisf\/ies formu\-la~\eqref{eq:phi_iso}. Using second formula
\eqref{eq:phir}, the formulae for the Fr\"olicher--Nijenhuis bracket of
two vector valued forms \cite[Appendix~A1]{grifone00}, and $[J, \mathbb{C}]=J$, we have
\begin{gather*}
3R=[J,\Phi]=[J, \lambda J + \eta\otimes \mathbb{C}] =
\left(d_J\lambda-\eta\right)\wedge J + d_J\eta \otimes \mathbb{C}. %\label{eq:3rle}
\end{gather*} Hence, the curvature tensor $R$ has the form
\eqref{eq:r_iso}.

We assume now that the curvature tensor $R$ has the form
\eqref{eq:r_iso}.  Using f\/irst formula \eqref{eq:phir} and the fact
that the inner product $i_S$ is a derivation of degree $-1$, we have
that the Jacobi endomorphism has the form
\begin{gather*}
\Phi=i_SR=i_S\alpha J + \left(i_S\beta-\alpha\right)\otimes
\mathbb{C}. %\label{eq:phiab}
\end{gather*} Hence, the spray $S$ is isotropic.
\end{proof}

We will use  Proposition \ref{prop:iso} and formula \eqref{eq:r_iso} in
Subsection \ref{subsec:pms} to show that isotropic sprays are
projectively metrizable sprays.


\subsection{Projectively related sprays} \label{subsec:prs}

Two sprays are projectively equivalent if their geodesics
coincide as oriented curves. Therefore, a spray is called projectively
metrizable if its geodesics coincide, as oriented curves, with the
geodesics of a Finsler space.

\looseness=-1
In \cite{bucataru09} it has been shown that the Helmholtz conditions for an
arbitrary semispray to be a~Lag\-rangian vector f\/ield can be
reformulated in terms of semi-basic $1$-forms. It has been shown also
that out of the four classic Helmholtz conditions only two of them are necessary and suf\/f\/icient
in the case of the projective metrizability problem for a spray. In this
subsection we obtain directly the two Helmholtz conditions, for
projective metrizability, in terms of semi-basic $1$-forms

\begin{Definition} \label{def:finsler}
By a \emph{Finsler function} we mean a continuous function $F: TM
\to \mathbb{R}$ satisfying the following conditions:
\begin{itemize}\itemsep=0pt
\item[i)] $F$ is smooth on $TM\setminus\{0\}$;
\item[ii)] $F$ is positive on $TM\setminus\{0\}$ and $F(x,0)=0$;
\item[iii)] $F$ is positively homogeneous of order $1$, which
means that $F(x,\lambda y)=\lambda F(x,y)$, for all $\lambda>0$
and $(x,y)\in TM$;
\item[iv)] the \emph{metric tensor} with components
\begin{gather*}
g_{ij}(x,y)=\frac{1}{2}\frac{\partial^2 F^2}{\partial y^i\partial
y^j} %\label{gij}
\end{gather*}
has rank $n$. \end{itemize}
\end{Definition}
According to Lovas \cite{lovas07}, conditions ii) and iv) of
Def\/inition~\ref{def:finsler}  imply that the metric tensor~$g_{ij}$ of
a Finsler function is positive def\/inite.

The regularity condition iv) of Def\/inition \ref{def:finsler} implies that the
Euler--Poincar\'e $2$-form of $F^2$, $\omega_{F^2}=dd_JF^2$, is non-degenerate
and hence it is a symplectic structure \cite{deleon89, morandi90}. Therefore, the equation
\begin{gather} i_Sdd_JF^2=-dF^2 \label{isddj}
\end{gather}
uniquely determine a vector f\/ield $S$ on $TM\setminus\{0\}$ that is called the
\emph{geodesic spray} of the Finsler function. Equation \eqref{isddj}
is equivalent to
\begin{gather} \mathcal{L}_Sd_JF^2=dF^2. \label{lsdj}
\end{gather}
Locally, the Euler--Poincar\'e $2$-form of $F^2$, $\omega_{F^2}=dd_JF^2$, can
be expressed as follows
\begin{gather*}
\omega_{F^2}=2g_{ij}\delta y^i \wedge dx^j. %\label{omegaf2}
\end{gather*}
\begin{Definition} \label{defn:finsler_metr}
A spray $S$ is called \emph{Finsler metrizable} if there exists a
Finsler function $F$ that satisf\/ies one of the two equivalent
conditions \eqref{isddj} or \eqref{lsdj}.
\end{Definition}
One can reformulate condition iv) of Def\/inition \ref{def:finsler}
in terms of the Hessian of the Finsler function $F$ as follows.
Consider
\begin{gather*}
h_{ij}(x,y)=F\frac{\partial^2 F}{\partial y^i\partial y^j}
%\label{hij}
\end{gather*} the \emph{angular metric} of the Finsler function.
The metric tensor $g_{ij}$ and the angular tensor $h_{ij}$ are
related by
\begin{gather*}
g_{ij}=h_{ij} + \frac{\partial F}{\partial y^i} \frac{\partial
F}{\partial y^j}. %\label{ghii}
\end{gather*} Metric tensor
$g_{ij}$ has rank $n$ if and only if angular tensor $h_{ij}$ has
rank $(n-1)$, see~\cite{matsumoto86}. Therefore, the regularity of
the Finsler function $F$ is equivalent with the fact that the
Euler--Poincar\'e $2$-form $\omega_{F}=dd_JF$ has rank $2n-2$.
\begin{Definition}\qquad
\begin{itemize} \itemsep=0pt \item[i)] Two sprays $S_1$ and $S_2$ are \emph{projectively equivalent}
if their geodesics coincide up to an orientation preserving
reparameterization. \item[ii)] A spray $S$ is \emph{projectively
metrizable} if it is projectively equivalent to the geodesic spray
of a Finsler function.
\end{itemize}
\end{Definition}
Two sprays $S_1$ and $S_2$ are projectively equivalent if and only
if there exists a $1$-homogeneous function $P\in
C^{\infty}(TM\setminus\{0\})$ such that $S_2=S_1-2P\mathbb{C}$,
\cite{aim93, shen01}.

Next theorem gives a characterization of projectively metrizable sprays in
terms of semi-basic $1$-forms on $TM\setminus\{0\}$.
\begin{Theorem} \label{thm:pm} A spray $S$ is projectively
metrizable if and only if there exists a semi-basic $1$-form
$\theta\in \Lambda^1_v(TM\setminus\{0\})$ such that
\begin{gather}
\operatorname{rank}\left(d\theta\right)=2n-2, \qquad i_S\theta>0,
\label{palgebric} \\ \mathcal{L}_{\mathbb{C}}\theta =0, \qquad
d_J\theta=0, \qquad d_h\theta=0. \label{cjh}\end{gather}
\end{Theorem}
\begin{proof}
We prove f\/irst that conditions \eqref{palgebric} and \eqref{cjh} are
necessary for the projective metri\-za\-bility problem of the spray $S$.
We assume that $S$ is projectively metrizable. Therefore, there
exists a~Finsler function $F$ with geodesic spray $S_F$ and a
$1$-homogeneous function $P$ on $TM\setminus\{0\}$ such that
$S=S_F-2P\mathbb{C}$. Consider $\theta=d_JF$, the Euler--Poincar\'e
$1$-form of the Finsler function~$F$. Due to the $1$-homogeneity
condition of $F$ it follows that $i_S\theta=\mathbb{C}(F)=F>0$. The non-degeneracy
of the Finsler function implies
$\operatorname{rank}\left(d\theta\right)=2n-2$. Since $\theta$ is
$0$-homogeneous it follows that
$\mathcal{L}_{\mathbb{C}}\theta=0$. Condition $d_J\theta=0$ is
also satisf\/ied since $d_J\theta=d_J^2F=0$.

\looseness=-1
It remains to show that $d_h\theta=0$. The geodesic spray $S_F$ is
uniquely determined by condi\-tion~\eqref{isddj}, from which it
follows that $S_F(F^2)=0$ and hence $S_F(F)=0$. Since $S_F$ also
satis\-f\/ies condition \eqref{lsdj} it follows that
$\mathcal{L}_{S_F}\left(F\theta\right)=FdF$, which implies
$\mathcal{L}_{S_F}\theta=dF$. Using $S=S_F-2P\mathbb{C}$ we obtain
that $\mathcal{L}_{S}\theta -2 \mathcal{L}_{P\mathbb{C}}\theta=dF$.
Using again the $0$-homogeneity of the semi-basic $1$-form $\theta$ it
follows $\mathcal{L}_{P\mathbb{C}}\theta= P
\mathcal{L}_{\mathbb{C}}\theta=0$ and hence
$\mathcal{L}_{S}\theta=dF$. We apply now $d_J$ to both sides of
this last relation and use the commutation rules
$\mathcal{L}_Sd_J-d_J\mathcal{L}_S=d_{[S,J]}=-d_{h}+d_{v}$ and
$dd_J+d_Jd=0.$ Therefore,
\[
-d_h\theta-d_v\theta=-dd_JF=d_JdF=d_J\mathcal{L}_{S}\theta=
\mathcal{L}_Sd_J\theta + d_{h}\theta- d_{v}\theta,
\] from where
it follows that $d_h\theta=0$.

We prove now that the conditions \eqref{palgebric} and \eqref{cjh}  are
suf\/f\/icient for the projective metrizability problem of the spray $S$. Consider $\theta\in \Lambda^1(TM\setminus\{0\})$
a semi-basic $1$-form that sa\-tis\-f\/ies conditions~\eqref{palgebric}
and~\eqref{cjh}. Def\/ine the function $F=i_S\theta$. Using the commutation rule
$i_Sd_J+d_Ji_S=\mathcal{L}_{\mathbb{C}}-i_{[S,J]}$ as well as
conditions $d_J\theta=0$ and $\mathcal{L}_{\mathbb{C}}\theta=0$ it
follows that $d_JF=d_Ji_S\theta=i_h\theta=\theta$. Hence $\theta$
is the Euler--Poincar\'e $1$-form of $F$. Now conditions~\eqref{palgebric} assure that $F$ is a Finsler function. Consider
the function $P\in C^{\infty}(TM\setminus\{0\})$ given by
$2P=S(F)/F$, which is $1$-homogeneous. We will show now that the
spray $\tilde{S}=S-2P\mathbb{C}$ satisf\/ies equation~\eqref{lsdj}
and hence it is the geodesic spray of the Finsler function~$F$.

Using the commutation rule $i_Sd_h+d_hi_S=\mathcal{L}_S-
i_{[S,h]}$ and the fact that $d_h\theta=0$ it follows
$0=i_Sd_h\theta=-d_hi_S\theta+ \mathcal{L}_S\theta-
i_{[S,h]}\theta$. Using the fact that $i_{[S,h]}\theta=dF\circ
J\circ \mathcal{L}_Sh=dF\circ v=d_vF$ it follows that
$\mathcal{L}_S\theta=d_hF+d_vF=dF$. We show now that $\tilde{S}$
satisf\/ies the same equation. Indeed
$\mathcal{L}_{\tilde{S}}\theta=
\mathcal{L}_{S-2P\mathbb{C}}\theta=dF$ since
$\mathcal{L}_{P\mathbb{C}}\theta=0$. From the def\/ining formula of
function $P$ is follows that
$\tilde{S}(F)=S(F)-2P\mathbb{C}(F)=S(F)-2PF=0$. Therefore
$\mathcal{L}_{\tilde{S}}d_JF^2=2F\mathcal{L}_{\tilde{S}}d_JF=2FdF=dF^2$
and hence $\tilde{S}$ is the geodesic spray of the Finsler
function~$F$.
\end{proof}

The second part of the proof of Theorem \ref{thm:pm} shows that if
there exists a semi-basic $1$-form $\theta$ on $TM\setminus\{0\}$ that
satisf\/ies the conditions \eqref{palgebric} and \eqref{cjh} then the
given spray $S$ is projectively related to the spray
\[
 S_F=S-\frac{\mathcal{L}_{S}(i_S\theta)}{i_S\theta}\mathbb{C},
 \]
  which is
the geodesic spray of the Finsler function $F=i_S\theta$. In this
case, the semi-basic $1$-form $\theta=\theta_i dx^i$ is the
Euler--Poincar\'e $1$-form of the Finsler function $F$,
$\theta=d_JF$. Therefore,
\begin{gather} \theta_i=\frac{\partial F}{\partial y^i}, \qquad
  h_{ij}=F\frac{\partial \theta_i}{\partial
    y^j}, \qquad Fd\theta=h_{ij} \delta y^i \wedge
  dx^j. \label{fdt} \end{gather}
Formulae \eqref{fdt} show the relation between a semi-basic $1$-form
$\theta$, a solution of the projective metrizability problem using
Theorem~\ref{thm:pm}, and the classic approach of the problem using the
multiplier matrix $h_{ij}$.


\section{Formal integrability for the projective metrizability problem}
\label{sec:fi}

Theorem \ref{thm:pm} provides necessary and suf\/f\/icient conditions for the projective
metrizability problem. These conditions consist of a set of algebraic equations~\eqref{palgebric},
and a set of dif\/ferential equations~\eqref{cjh}. In this section,
we study the set of dif\/ferential equations~\eqref{cjh} using Spencer's
technique of formal integrability \cite{BCG91, grifone00}  and following some of the techniques
used for studying the Finsler metrizability problem, which were developed in~\cite{muzsnay06}.


\subsection{Formal integrability} \label{subsec:fi}
In this subsection, we recall f\/irst the basic notions of formal integrability
\cite{BCG91, grifone00} and then we apply it to the system \eqref{cjh}.

Consider $E$ a vector bundle over the base manifold $M$. For a
section $s$ of $E$ and $k\geq 1$ we denote by $j^k_xs$ the $k$th
order jet of $s$ at the base point~$x$ in~$M$. The bundle of $k$th
order jets of sections of $E$ is denoted by~$J^kE$. For two vector
bundles~$E$ and~$F$ over the same base manifold~$M$, a linear
partial dif\/ferential operator of order~$k$,
\[
P: \  \operatorname{Sec}(E) \to \operatorname{Sec}(F),
\]  can be identif\/ied with a morphism of vector
bundles over $M$, $p^{0}(P): J^{k}E \to
F$. We will also consider the $l$th order jet prolongation of the dif\/ferential
operator $P$, which will be identif\/ied with the morphisms of vector bundles
over $M$, $p^{l}(P): J^{k+l}E \to J^lF$, def\/ined by
\[
p^{l}(P)\big(j^{k+l}_xs\big)=j^l_x(Ps).
\] We will denote by
$R^{k+l}_x(P)=\operatorname{Ker} p_x^{l}(P)\subset J^{k+l}_xE$ the
space of $(k+l)$th order formal solutions of $P$ at $x$ in $M$.
\begin{Definition} \label{def:fi} The dif\/ferential operator $P$ is called \emph{formally integrable} at
$x$ in $M$ if $R^{k+l}(P)$ is a vector bundle over $M$, for all $l\geq 0$,
and the map $\bar{\pi}^{k+l-1}_x: R^{k+l}_x(P) \to R^{k+l-1}_x(P)$ is onto for all
$l\geq 1$. \end{Definition}
In the analytic case, formal integrability implies
existence of analytic solutions for arbitrary initial data, see
\cite[p.~397]{BCG91}.

Denote by $\sigma^k(P): S^k(M)\otimes E \to F$ the symbol of $P$,
which is def\/ined by the highest order terms of the dif\/ferential operator
$P$, and by $\sigma^{k+l}(P): S^{k+l}(M)\otimes E \to S^l(M)
\otimes F$ the symbol of the $l$th order prolongation of $P$. For
each $x$ in $M$, we write
\begin{gather*}
g^k_x(P) = \operatorname{Ker}\sigma^k_x(P),
\\ g^k_x(P)_{e_1\dots e_j}  =  \{A \in g^k_x(P)| i_{e_1}A=\cdots =
i_{e_j}A=0\},\qquad j\in \{1,\dots ,n\},
\end{gather*}
where  $\{e_1,\dots ,e_n\}$ is a basis of $T_xM$. Such a basis is
called \emph{quasi-regular} if it satisf\/ies
\begin{gather} \dim g^{k+1}_x(P)= \dim g^k_x(P) + \sum_{j=1}^n \dim
g^k_x(P)_{e_1\dots e_j}. \label{quasir}\end{gather} The symbol
$\sigma^k(P)$ is called \emph{involutive} at
$x$ in $M$ if there exists a quasi-regular basis of $T_xM$.

In this work we will address the projective metrizability problem by
discussing f\/irst the formal integrability of the system \eqref{cjh}. For
this we will use the two suf\/f\/icient conditions provided by
Cartan--K\"ahler theorem.

\medskip

\noindent {\bf Theorem} [Cartan--K\"ahler]. {\it Let $P$ be a linear partial
differential operator of order $k$. Suppose $g^{k+1}(P)$ is a
vector bundle over $R^k(P)$. If the map $\overline{\pi}^k:
R^{k+1}(P) \to R^k(P)$ is onto and the symbol $\sigma^k(P)$ is involutive, then
$P$ is formally integrable.}

\medskip

In order to study the formal integrability of the system
\eqref{cjh} we consider the f\/irst-order partial dif\/ferential operator $P_1:
\Lambda^1_v(TM\setminus\{0\}) \to \Lambda^1_v(TM\setminus\{0\})
\oplus \Lambda^2_v(TM\setminus\{0\}) \oplus
\Lambda^2_v(TM\setminus\{0\})$, which we call \emph{the projective
  metrizability operator}
\begin{gather}
\label{p1} P_1  = \left({\mathcal L}_{\mathbb C}, d_J, d_h
\right).
\end{gather}
Since $\mathbb{C} $ and $J$ are vector valued, semi-basic $0$ and
respectively $1$-forms and
$h$ is a vector valued almost semi-basic $1$-form,  according to Lemma \ref{lem:dlomega}, all dif\/ferential
operators $\mathcal{L}_{\mathbb{C}}$, $d_J$, $d_h$ preserve
semi-basic forms. Therefore, the dif\/ferential operator $P_1$ is well
def\/ined.

